\section{Introduction}
\label{sec:intro}

A neutron star maps an equation of state into a collection of exterior and bulk observables through the relativistic stellar-structure equations. Individual observables such as the radius, moment of inertia, and tidal deformability retain substantial dependence on the EOS. Certain combinations, however, vary much less. The I--Love--Q relations provide the canonical example, relating the dimensionless moment of inertia, tidal deformability, and spin-induced quadrupole moment with weak EOS dependence \cite{yagi2013iloveq}. Their accuracy has motivated both practical applications and a long-running question: why does eliminating the stellar-configuration variable remove so much of the EOS dependence?

Several complementary explanations already exist and set the starting point for this work. Approximate self-similarity of isodensity contours and the dominance of the outer core were identified as important ingredients \cite{yagi2014why}. The proximity of realistic compact stars to stiff or incompressible configurations provides another perspective \cite{sham2015unveiling}. Perturbative calculations sharpened this statement: the I--Love relation can be stationary with respect to EOS variations around the incompressible limit \cite{chan2016stationarity}, and in Newtonian gravity the stationarity of moment--Love relations extends to arbitrary density perturbations about an incompressible star \cite{yip2017tidal}. More recently, Hu, Gao, and Shao separated linear deviations of the I--C and I--Love relations into an EOS-difference factor and a background-structure factor \cite{hu2026linear}, while Kamata, Minamiguchi, and Minato derived asymptotic universal relations directly from the relativistic stellar equations and boundary conditions \cite{kamata2026asymptotic}. These results explain or quantify why already identified relations are insensitive in controlled limits and provide important first-order and asymptotic antecedents for the construction below.

The modern EOS literature also makes clear that the term universal is conditional. Flexible nonparametric EOS ensembles produce larger deviations than compact catalogs of named EOSs and expose relations whose apparent accuracy depends on the EOS prior \cite{legred2024assessing}. Statistical and machine-learning methods have been used to search directly for combinations of observables with small EOS scatter \cite{manoharan2024finding,kruger2026sbi}. Automatic differentiation through the stellar-structure problem is practical \cite{soma2023ad}, while sensitivity matrices have been used to determine which microscopic EOS parameters most strongly affect neutron-star observables \cite{cruzcamacho2026sensitivity}.

These developments suggest a different formulation. Instead of beginning with a proposed relation and measuring its scatter, we begin with the differential map from EOS function space to observable space. We ask which normal directions in observable space are weakly excited by physically weighted EOS perturbations after displacement along the stellar sequence has been removed. This turns quasi-universality into a local geometric statement and separates three questions that are usually combined:
\begin{enumerate}
    \item Does the EOS response possess a small transverse direction locally?
    \item Do those local directions integrate to a global relation?
    \item What controls the finite departure from that relation?
\end{enumerate}

The distinction matters. A small local response need not define a global function \(F(\yvec)\), and a first-order stationary relation can still acquire finite scatter from second-order curvature. Conversely, a breakdown of universality may be attributable to a specific density-localized EOS direction rather than to an undifferentiated increase in catalog scatter.

This paper develops that construction for cold, barotropic neutron stars in general relativity. We initially restrict the observable space to
\begin{equation}
    \yvec =
    \left(\ln C,\ln\Ibar,\ln\Lamtwo\right),
    \qquad
    C\equiv M/R,\qquad
    \Ibar\equiv I/M^3,
\end{equation}
and treat the known I--Love relation as a stringent validation target. Compactness--Love provides a deliberately weaker control. We do not assume either relation when constructing the response geometry.

The calculation recovers a strongly separated soft normal mode without using an empirical I--Love fit, and that mode remains nearly parallel to the I--Love normal over broad causal EOS deformations.  The resulting one-form is approximately integrable, while second-order curvature controls the leading finite departure from the local soft surface.  A controlled localized sound-speed softening experiment then shows that the I--Love projection can degrade substantially while the full low-response hierarchy survives through a rotation of the soft covector.  The extra sensitivity can be localized to the responsible stellar layer.

Section~\ref{sec:geometry} defines the EOS response operator, the induced metric in observable space, the stellar-sequence quotient, integrability, and the second-order curvature terms. Section~\ref{sec:stars} specifies the stellar-structure problem and EOS perturbations. Section~\ref{sec:numerics} gives the numerical validation. Section~\ref{sec:localresults} presents the local response hierarchy and full normal spectrum, Sec.~\ref{sec:global} tests integrability and reconstructs the scalar potential, Sec.~\ref{sec:curvature} tests the finite-radius curvature expansion, and Sec.~\ref{sec:breakdown} isolates a controlled breakdown mechanism.
